\section{Урок 6: базовые фигуры и миссия по точкам}

\subsection{Цели урока}
\begin{itemize}
\item Сформировать миссию из массива точек.
\item Понять параметризацию окружности через угол и радиус.
\item Совместить движение и видеонаблюдение.
\end{itemize}

\subsection{Опорные листинги}
\lstinputlisting[style=GeoskanStyle,language=pythoncustom,caption={Маршрут по массиву точек с кадрами камеры}]{../code/python/examples/mission_examples/mission_with_frames.py}
\lstinputlisting[style=GeoskanStyle,language=pythoncustom,caption={Окружность по дискретным точкам}]{../code/python/examples/mission_examples/circle_flight.py}

\subsection{Формула окружности}
\[
x = r\cos\varphi,\quad y = r\sin\varphi
\]
где $r$ — радиус, $\varphi$ — текущий угол в радианах.

\begin{table}[H]
\centering
\begin{tabular}{cccc}
\toprule
$\varphi$ & $x/r$ & $y/r$ & Интерпретация \\
\midrule
0 & 1.00 & 0.00 & старт на оси $x$ \\
$\pi/2$ & 0.00 & 1.00 & верхняя точка \\
$\pi$ & -1.00 & 0.00 & противоположная точка \\
$3\pi/2$ & 0.00 & -1.00 & нижняя точка \\
\bottomrule
\end{tabular}
\caption{Опорные значения для генерации фигуры}
\end{table}

\begin{taskbox}[Минимальный набор упражнений]
\begin{enumerate}
\item Увеличьте \texttt{number\_of\_points} до 16 и сравните плавность траектории.
\item Постройте квадратный маршрут из массива \texttt{waypoints}.
\item Добавьте возврат в стартовую точку на высоте 1 м перед посадкой.
\end{enumerate}
\end{taskbox}

\subsection{Критерии успешности}
\begin{enumerate}
\item Фигура замыкается без резкого рывка на последнем участке.
\item Камера непрерывно показывает кадры в ходе миссии.
\item Все этапы завершены с безопасной посадкой.
\end{enumerate}
